\chapterbanner{第3章}{\jas{燃料・エネルギー}}{問8〜問14 \jas{・} 過去失分率 49\%}
\label{ch03-start}
\begin{keybox}{本章用法}
統計題：先圈\textbf{年度}與\textbf{分母}（\ja{一次エネルギー供給} vs 発電電力量）。「無限」「\ja{必ず}」「過去最短」等絕對化用語先懷疑；正解多為「長所＋課題」並列的敘述。數值標 ※ 者每年變動，考前用最新『\ja{エネルギー白書}』確認。

\end{keybox}
\secthead{3.1 \jas{エネルギー統計・政策}}{\#212\,\#199\,\#283\,\#264\,\#109\,\#228\,\#241\,\#242\,\#5\,\#143\,\#74H27}
\subsecthead{\jas{一次エネルギーと二次エネルギー}}{\#5\,\#143}
判定：\ja{自然界からそのまま採ったか}（一次）、\ja{工場で転換・加工したか}（二次）。

\begin{gtable}{YY}
\rowcolor{thead}\textbf{一次} & \textbf{二次} \\
{}\ja{原油・石炭・天然ガス・ウラン} & 電力（\ja{最も典型}） \\
{}\ja{太陽・風力・水力・地熱・バイオマス} & \ja{ガソリン・灯油・軽油・重油・ナフサ・}LPG\ja{・都市ガス} \\
{}— & \ja{コークス・コールタール・石炭ガス}；\hl{水素}\ja{・バイオエタノール} \\
\end{gtable}
「一次」\ja{は取り出す順番の意味で}、\ja{再生可能かどうかとは無関係}。

\subsecthead{\jas{日本の数字}（※ \jas{最新値を確認}）}{\#199\,\#283\,\#212\,\#74H27}
\begin{gtable}{YY}
\rowcolor{thead}\textbf{指標} & \textbf{概略値} \\
{}\ja{一次エネルギー供給}（2017 年度） & 石油 ≈39\%（最大）＞ 石炭 ≈25\% ≈ \ja{天然ガス} ≈23\% ＞ \ja{再エネ}（\ja{水力含む}）≈11\% ＞ 原子力 ≈1.4\% \\
{}発受電電力量（2017 年度） & \hl{LNG ≈40\%（最大）}、石炭 ≈32\%、\ja{再エネ} ≈16\%、石油 ≈9\%、原子力 ≈3\% \\
{}発電（2023 年度概略） & LNG\ja{・石炭} 各 ≈3 割、\ja{再エネ} 2 割強（太陽光 ≈1 \ja{割で最大}）、原子力 ＜1 割 \\
{}\ja{最終エネルギー消費} ÷ 一次供給 & \hl{約 2/3（≈65\%）}；\ja{損失の主因は火力発電の廃熱}（効率 ≈40\%） \\
{}\ja{エネルギー自給率} & ≈12〜15\%（1960 年 ≈58\%） \\
{}\ja{原油の中東依存} & 9 割以上；LNG \ja{はオーストラリア} ≈4 \ja{割が最大} \\
{}\ja{石油の一次供給比率} & 1973 年 ≈75\% → 2020 年度 ≈36\%（1965 年度以来最低） \\
\end{gtable}
\begin{itemize}
\item{} \ja{同じ石炭でも一次} ≈25\%、発電 ≈32\%：\textbf{\ja{分母を確認}}。\ja{石油は運輸・石化原料用なので発電では小さく}、LNG \ja{はほぼ発電用なので発電で最大}。
\item{} 2014 \ja{年度は原子力} 0\%（特殊年度）。
\end{itemize}
\subsecthead{\jas{化石燃料・可採年数}}{\#242\,\#264}
R/P ＝ 確認埋蔵量 ÷ 年生産量。\ja{埋蔵量は技術・価格で上方修正されるので}\textbf{\ja{カウントダウンではない}}。石油 ≈50 年（1980 \ja{年代は} ≈30 年、「過去最短」\batsu{}）、\ja{天然ガス} ≈50 年、石炭 ＞100 年（\ja{分布が分散}、CO\textsubscript{2} 最多）。資源量（\ja{技術的に可採}）≠ 埋蔵量（\ja{経済的に可採}）。\ja{シェールガスの技術的可採資源量} ＞ \ja{在来型天然ガスの確認埋蔵量}。

\subsecthead{気候政策（R4-Q8 失分）}{\#228\,\#241\,\#264}
\begin{gtable}{lYY}
\rowcolor{thead}\textbf{} & \textbf{京都議定書（1997）} & \textbf{\jas{パリ協定}（2015 \jas{採択・}2016 発効）} \\
{}\ja{目標の決め方} & \ja{トップダウン}（日本 −6\%） & \ja{ボトムアップ}（\ja{各国が} NDC \ja{を自主設定}） \\
{}法的性格 & 達成義務 & \ja{提出・更新は義務}、\textbf{\ja{達成は義務でない}} \\
{}対象 & 先進国 & 全締約国 \\
\end{gtable}
\begin{itemize}
\item{} \ja{パリ協定}：工業化前比 2 ℃ \ja{より十分低く}、1.5 ℃ \ja{に抑える努力}（「5 ℃」\batsu{}）。
\item{} 日本：2050 \ja{年カーボンニュートラル}（2020 宣言）、2030 年度 −46\%（2013 年度比、\ja{自主的に表明}）。\ja{温室効果ガスの} ≈9 \ja{割が} CO\textsubscript{2}；\ja{直接排出は}\textbf{\ja{エネルギー転換部門}}\ja{が最大}、\ja{電熱配分後は産業部門が最大}。
\item{} \ja{代替フロン} HFC：\ja{オゾン層は壊さないが温暖化係数が高く排出増加中}（\ja{キガリ改正}）。
\end{itemize}
\subsecthead{\jas{再生可能エネルギー}（R4-Q13 失分）}{\#109\,\#264}
\begin{itemize}
\item{} 風力 P ＝ ½ρAv\textsuperscript{3}（A ∝ D\textsuperscript{2}）→ \ja{大型化・洋上化}。\ja{ダム式水力は短期天候に左右されず}、\ja{揚水は調整電源}（\ja{巨大な電池}）。
\item{} 地熱：\ja{日本は資源量世界第} 3 \ja{位だが設備容量は伸びない}（\ja{国立公園・温泉との調整}）。
\item{} \ja{バイオマス比率は}\textbf{途上国 ≫ 先進国}（\ja{伝統的バイオマス}＝\ja{薪・炭}；H29-Q8）。「\ja{無限に手に入る}」\batsu{}。
\item{} 2015 年：太陽光 中国 ＞ \ja{ドイツ} ＞ 日本、風力 中国 ＞ 米国 ＞ \ja{ドイツ}（2 \ja{位が違う}）。
\end{itemize}
\secthead{3.2 石油精製}{\#243\,\#111\,\#145\,\#76H27\,\#77\,\#262\,\#263\,\#201\,\#7\,\#42}
\begin{figblock}
\begin{tikzpicture}[font=\scriptsize,node distance=4mm,
  b/.style={draw,rounded corners=1pt,align=center,inner sep=2pt,fill=thead},
  p/.style={align=left,inner sep=1pt}]
\node[b] (crude) {原油};
\node[b,right=6mm of crude] (atm) {常圧蒸留};
\node[p,right=5mm of atm,yshift=9mm] (lpg) {LPG};
\node[p,right=5mm of atm,yshift=4.5mm] (ln) {軽質ナフサ → 異性化・熱分解（エチレン）};
\node[p,right=5mm of atm,yshift=0mm] (hn) {重質ナフサ → 接触改質};
\node[p,right=5mm of atm,yshift=-4.5mm] (ker) {灯油・軽油 → 水素化精製};
\node[b,below=7mm of atm] (vac) {減圧蒸留};
\node[p,right=5mm of vac,yshift=3mm] (vgo) {減圧軽油 → 脱硫 → FCC・水素化分解};
\node[p,right=5mm of vac,yshift=-1mm] (lub) {潤滑油留分 → 溶剤精製・脱ろう};
\node[p,right=5mm of vac,yshift=-5mm] (res) {減圧残油 → アスファルト・重油};
\draw[->] (crude)--(atm);
\foreach \n in {lpg,ln,hn,ker} \draw[->] (atm.east)--(\n.west);
\draw[->] (atm.south)--node[left]{常圧残油}(vac.north);
\foreach \n in {vgo,lub,res} \draw[->] (vac.east)--(\n.west);
\end{tikzpicture}
\end{figblock}
\subsecthead{\jas{装置の比較}（必考）}{\#243\,\#111\,\#145\,\#76H27\,\#77}
\begin{gtable}{>{\hsize=0.8\hsize}Y>{\hsize=0.9\hsize}Y>{\hsize=0.9\hsize}Y>{\hsize=1.4\hsize}Y}
\rowcolor{thead}\textbf{装置} & \textbf{原料} & \textbf{触媒} & \textbf{要點} \\
{}流動接触分解 FCC & 脱硫減圧軽油 & \hl{\ja{ゼオライト}}（固体酸） & \ja{カルボカチオン機構} → \ja{分枝化・芳香族化} → \textbf{\ja{高オクタン価}}；C3\ja{・}C4 \ja{オレフィン副生}；\ja{再生器でコークを燃やした熱で反応熱を賄う}（加熱炉不要） \\
{}接触改質（\ja{リフォーマー}） & \hl{重質}\ja{ナフサ}（C7〜C10） & \hl{Pt}（Pt–Re、Pt–Sn）／\ja{塩素化アルミナ}（二元機能） & \ja{ナフテン脱水素・脱水素環化} → 芳香族＋\textbf{H\textsubscript{2} 副生}（\ja{製油所最大の水素源}）；吸熱 → 中間加熱炉 \\
{}\ja{アルキル化} & \ja{イソブタン}＋\ja{ブテン・プロピレン} & H\textsubscript{2}SO\textsubscript{4} \ja{または} HF（\hl{\ja{塩酸は使わない}}） & \ja{イソオクタン主体のアルキレート}：\ja{高オクタン価・芳香族}/\ja{硫黄なし}；\ja{結合なので発熱} \\
{}異性化 & \ja{軽質ナフサ}（C5〜C6） & — & 直鎖 → \ja{分枝でオクタン価}↑ \\
{}水素化脱硫 HDS & 各留分 & Co–Mo\ja{・}Ni–Mo/Al\textsubscript{2}O\textsubscript{3} & S → H\textsubscript{2}S（\ja{クラウス法で硫黄回収}） \\
{}水素化分解 & 重質油 & \ja{ゼオライト}＋Ni/W & 分解＋水添、H\textsubscript{2} \ja{を消費} \\
\end{gtable}
\begin{trapbox}{\jas{定番の入れ替え}（H29-Q10\jas{・}R3-Q10\jas{・}R7-Q9）}
「FCC \ja{の活性成分は白金}」\batsu{}（\ja{白金は改質}；FCC \ja{は再生で} 700 ℃ \ja{の空気燃焼を繰り返し}、\ja{硫黄・金属も含むので} Pt \ja{は使えない}）。「\ja{分解するとオクタン価が下がる}」\ja{は熱分解の話で} FCC \ja{では} \batsu{}。\ja{改質の原料は}\textbf{重質}\ja{ナフサ}（軽質 \batsu{}）。

\end{trapbox}
\subsecthead{\jas{接触改質の操作方向}}{\#145\,\#76H27}
\begin{gtable}{cYYl}
\rowcolor{thead}\textbf{操作} & \textbf{\jas{オクタン価}} & \textbf{液収率} & \textbf{\jas{コーク・触媒寿命}} \\
{}温度↑ & ↑ & ↓（\ja{水素化分解でガス化}） & 悪化 \\
{}圧力↓ & ↑（1 mol → 4 mol、\ja{ルシャトリエ}） & ↑ & 悪化 \\
\end{gtable}
\ja{圧力を下げると両方良くなるが}、\ja{コークが増える} → \ja{低圧の} \textbf{CCR（連続触媒再生式）}\ja{が主流}。\ja{水素分圧を上げるとコークは減る}。

\subsecthead{\jas{脱硫・芳香族・減圧蒸留}}{\#262\,\#263\,\#201\,\#42}
\begin{itemize}
\item{} HDS \ja{の難易}：\ja{メルカプタン} ＜ \ja{スルフィド} ＜ \ja{チオフェン} ＜ \ja{ジベンゾチオフェン}（4,6-\ja{ジメチル体が最難}）；\ja{高沸点ほど難しい}。\ja{常圧残油の直接脱硫は} Ni\ja{・}V \ja{で触媒が被毒し過酷}、間接脱硫（\ja{減圧軽油のみ}）\ja{は穏やか}。「\ja{除去しやすい順}」＝\ja{最初が最も易しい}。
\item{} BTX：\ja{重質ナフサ改質のリフォーメート}（＋\ja{ナフサ熱分解の分解ガソリン}）→ \textbf{\ja{スルホラン}}抽出（高極性；CCl\textsubscript{4} \ja{は無極性で不可}）→ 蒸留。p-\ja{キシレンは} m-\ja{体と沸点差} ≈1 ℃ → 擬似移動床吸着（\ja{パレックス法}）。
\item{} \ja{常圧蒸留で得られるのは} LPG\ja{・ナフサ・灯油・軽油・常圧残油まで}。\ja{潤滑油は常圧で加熱すると先に熱分解} → \textbf{減圧蒸留}\ja{が必要}（H24-Q9\ja{・}R2-Q10）。
\end{itemize}
\secthead{3.3 \jas{石油製品の規格・物性}}{\#112\,\#110\,\#75H27\,框8\,\#76\,\#43\,\#144\,\#78\,\#284\,\#78H27\,\#146\,\#8\,\#200\,\#172\,\#6}
\subsecthead{\jas{製品と品質指標}}{\#112\,框8\,\#43}
\begin{gtable}{lYYll}
\rowcolor{thead}\textbf{製品} & \textbf{指標} & \textbf{定義} & \textbf{有利} & \textbf{不利} \\
{}\ja{ガソリン} & \ja{オクタン価}（\ja{耐ノック}） & \ja{イソオクタン} 100、n-\ja{ヘプタン} 0 & \ja{芳香族・分枝} & \ja{直鎖パラフィン} \\
{}灯油 & \hl{煙点}（mm、\ja{高いほど良い}） & \ja{煙を出さずに燃える最大炎高さ} & \ja{パラフィン} & 芳香族 \\
{}軽油 & \ja{セタン価}（着火性） & n-\ja{セタン} 100、α-\ja{メチルナフタレン} 0 & \ja{直鎖パラフィン} & 芳香族 \\
{}\ja{ジェット燃料} & \ja{煙点・析出点} & — & \ja{パラフィン} & 芳香族 \\
{}重油 & 動粘度（A\ja{・}B\ja{・}C 重油） & — & — & — \\
\end{gtable}
\begin{memobox}{\jas{一本の線}}
\ja{芳香族が有利なのは}\textbf{\ja{ガソリンだけ}}。灯油（煙点↓）\ja{・軽油}（\ja{セタン価}↓）\ja{・ジェット燃料では不利}。\ja{オクタン価とセタン価は逆相関}。\ja{根はどれも} H/C 比。

\end{memobox}
\subsecthead{\jas{物性の向き}}{\#144\,\#110\,\#75H27\,\#76}
\begin{gtable}{Ycll}
\rowcolor{thead}\textbf{性質} & \textbf{炭素数↑} & \textbf{分枝↑} & \textbf{芳香族性↑} \\
{}\ja{沸点・粘度} & ↑ & ↓（接触面積↓） & ↑ \\
{}密度 & ↑ & \ja{やや}↓ & ↑ \\
{}\ja{重量当たり発熱量} & \ja{やや}↓ & \ja{ほぼ不変} & ↓（H/C 小） \\
{}\ja{体積当たり発熱量} & ↑ & — & ↑ \\
{}\ja{オクタン価} & ↓ & ↑ & ↑ \\
{}\ja{セタン価} & ↑ & ↓ & ↓ \\
\end{gtable}
\begin{itemize}
\item{} \ja{重量発熱量は} \textbf{H/C 比}\ja{で決まる}（H ≈142、C ≈33 MJ/kg）：\ja{ガソリン} ＞ 灯油 ＞ 軽油 ＞ 重油。\ja{体積基準は}\textbf{逆順}（\ja{密度が効く}）。\ja{ベンゼンの重量発熱量は} n-\ja{ヘキサンより}\textbf{\ja{小さい}}（H29-Q9）、\ja{体積では大きい}。
\item{} 高位発熱量（\ja{水蒸気が凝縮}、\ja{潜熱を含む}）＞ 低位発熱量。H 1 kg → H\textsubscript{2}O 9 kg。\ja{石油製品で差} 6〜7\%、\ja{天然ガス} ≈10\%。\ja{ボンベ熱量計は高位}。\ja{熱効率は通常低位で計算}。
\item{} \ja{比熱は組成によらずほぼ一定}（≈2 kJ/(kg·K)）、\ja{温度とともに増加}。
\item{} 動粘度 ν ＝ 絶対粘度 μ ÷ 密度 ρ（m\textsuperscript{2}/s；Pa·s ÷ kg/m\textsuperscript{3} \ja{で次元検算}）。\ja{日本語の}「動粘度」＝kinematic。
\item{} 低温：\textbf{\ja{曇り点} ＞ \ja{目詰まり点}（CFPP）＞ 流動点}。「\ja{流動する最低温度を曇り点という}」\batsu{}（\ja{定義のすり替え}）。
\item{} 沸点（同炭素数）：芳香族 ＞ 環状 ＞ 直鎖 ＞ 分枝。\ja{融点だけは高対称分枝で逆転あり}（\ja{ネオペンタン}）。
\end{itemize}
\subsecthead{\jas{規格の数字}}{\#78\,\#284\,\#78H27\,\#146\,\#200\,\#8}
\begin{gtable}{YYY}
\rowcolor{thead}\textbf{製品} & \textbf{\jas{炭素数・沸点}} & \textbf{規格} \\
{}LPG & C3〜C4 & \ja{プロパン}（臨界温度 97 ℃）\ja{は常温で加圧だけで液化}；LNG（\ja{メタン} −82 ℃）\ja{は} −162 ℃ 冷却 \\
{}\ja{ガソリン}（JIS K 2202） & C4〜C10、≈30〜200 ℃ & \ja{ハイオク} 96 \ja{以上・レギュラー} 89 以上；\ja{ベンゼン} \hl{1 vol\% 以下}；硫黄 10 ppm 以下 \\
{}灯油 & C11〜C13（≈C9〜C18）、170〜250 ℃ & 引火点 40 ℃ 以上；1 \ja{号の煙点} 23 mm 以上；燃料電池用硫黄 ≤10 ppm \\
{}軽油（JIS K 2204） & C14〜C20 & \ja{流動点で号数}：特1号 +5 → 3号 −20 → 特3号 −30 ℃ \\
{}\ja{ジェット燃料} & — & 民間機＝\hl{灯油型}（Jet A-1）、軍用＝広沸点範囲型（JP-4） \\
{}重油 & C20 以上 & A \ja{重油は加熱不要}、C \ja{重油は加熱必須}（250〜1000 mm\textsuperscript{2}/s） \\
\end{gtable}
\begin{itemize}
\item{} \ja{ガソリンの蒸留性状}：前端（T10）\ja{高すぎ} → 低温始動難、\ja{低すぎ} → \ja{ベーパーロック}；中間（T50）→ \ja{加速・暖機}；後端（T90）\ja{高すぎ} → \ja{オイル希釈}。
\item{} 消防法第 4 類：特殊引火物（\ja{ジエチルエーテル・}CS\textsubscript{2}）、第 1 石油類 ＜21 ℃（\ja{ガソリン}）、第 2 類 21〜70 ℃（\ja{灯油・軽油}）、第 3 類 70〜200 ℃（重油）、第 4 類（潤滑油）。
\item{} \textbf{\ja{引火点と発火点は逆向き}}：\ja{ガソリンは引火点} −40 ℃ \ja{以下だが発火点} ≈300 ℃；\ja{灯油は引火点} 40 ℃ \ja{以上・発火点} ≈220 ℃（\ja{長鎖ほど自己着火しやすい}＝\ja{セタン価と同原理}）。\ja{重油も水より軽い} → \ja{水で消火しない}。CS\textsubscript{2} \ja{の水中保存は発火点が極低いから}。
\item{} \ja{ジェット燃料は直留}＋水素化精製（\ja{芳香族・オレフィンを避ける}）、\ja{ガソリンは} FCC\ja{・改質で芳香族を}\textbf{\ja{増やす}}（\ja{逆の思想}）。
\end{itemize}
\secthead{3.4 潤滑油}{\#113\,\#265\,\#287\,\#173\,\#244\,\#79\,\#79H27\,\#44\,\#9}
\begin{gtable}{YYY}
\rowcolor{thead}\textbf{性質} & \textbf{\jas{パラフィン系}} & \textbf{\jas{ナフテン系}} \\
{}粘度指数 & \hl{高} & 低 \\
{}流動点 & 高（\ja{ワックス}） & 低 \\
{}溶解力（\ja{アニリン点}） & 低（\ja{アニリン点高}） & 高 \\
{}用途 & \ja{エンジン油・ギヤ油・油圧作動油} & \hl{\ja{冷凍機油・電気絶縁油}}\ja{・金属加工油} \\
\end{gtable}
\begin{trapbox}{\jas{用途の入れ替え}（H29-Q12\jas{・}R3-Q12\jas{・}R5-Q12）}
「\ja{冷凍機油・電気絶縁油はパラフィン系}」\batsu{}。\ja{低温用途にワックスは禁物}。HFC \ja{冷媒用はエステル系}（POE）合成油。高 VI \ja{基油は水素化分解・水素化異性化で作る}（\ja{接触改質は芳香族を増やすので逆}）。

\end{trapbox}
\begin{itemize}
\item{} 表示：\textbf{SAE}＝粘度（5W-30：W＝Winter、5W \ja{は低温側}、30 \ja{は} 100 ℃ 側）；\textbf{API}＝品質（S＝\ja{ガソリン用}、C＝\ja{ディーゼル用}、\ja{後ろの文字ほど新しい}：SL → SN → SP）。
\item{} 用途：ATF（有段 AT、CVT \ja{は} CVTF、\ja{互換なし}）、\ja{タービン油}（\ja{酸化安定性が最重要}）、\ja{ドライサンプ}＝大型船用低速 2 \ja{サイクルディーゼル・レース車}。
\item{} 添加剤：清浄分散剤（Ca\ja{・}Mg \ja{スルホネート}）、酸化防止剤（\ja{ヒンダードフェノール・}ZnDTP）、粘度指数向上剤（PMA\ja{・}OCP：\ja{高温で糸まりが広がる}；\ja{高分子量ほど増粘効果大だがせん断安定性は}\textbf{\ja{悪い}}）、流動点降下剤（\ja{ワックス結晶に吸着}）、極圧剤（S\ja{・}P\ja{・}Cl）。配合量：自動車用 5〜25\%、工業用 ≤5\%。
\item{} \ja{エンジン油の五作用}：\ja{潤滑・冷却・清浄分散・密封・防錆}。\ja{冷却の対象はピストン・軸受}（\hl{\ja{オイルフィルターではない}}；H28-Q12＝R2-Q12）。\ja{全損式マシン油は安価な油}。
\end{itemize}
\secthead{3.5 石炭}{\#114\,\#147\,\#80\,\#80H27\,\#10\,\#245}
石炭化度：泥炭 → 褐炭 → 亜瀝青炭 → 瀝青炭 → 無煙炭。

\begin{gtable}{YYY}
\rowcolor{thead}\textbf{\jas{単調に増加}} & \textbf{\jas{単調に減少}} & \textbf{\jas{単調でない}（\jas{瀝青炭で最大}）} \\
{}\ja{固定炭素・燃料比・密度・}\hl{着火点} & \ja{揮発分・水分・酸素・}H/C & \hl{発熱量}\ja{・}\hl{粘結性} \\
\end{gtable}
\begin{itemize}
\item{} 燃料比 ＝ 固定炭素 ÷ 揮発分。工業分析：水分（107 ℃）\ja{・揮発分}（900 ℃ 空気遮断）\ja{・灰分}（815 ℃ 燃焼）\ja{を測り}、\ja{固定炭素は}\textbf{\ja{差引き}}（\ja{実測しない}）。
\item{} \ja{褐炭は揮発分が多い} → \textbf{\ja{着火点が低く}}\ja{自然発火しやすい}（「\ja{褐炭の着火点は瀝青炭より高い}」\batsu{}；H29-Q13\ja{・}R3-Q13）。\ja{無煙炭は青く短い炎・無煙}。
\item{} 粘結性（約 400 ℃ \ja{で軟化溶融} → 再固化）\ja{は瀝青炭だけ} →\textbf{原料炭}（\ja{コークス用}、製鉄）；\ja{燃やすだけの石炭}＝\textbf{一般炭}。\ja{名称の入れ替えに注意}。\ja{コークスの高炉での役割}＝\ja{還元剤・熱源・通気}。
\item{} 乾留（空気遮断加熱）\ja{の生成物}：\ja{コークス}（固体）、\ja{石炭ガス}（H\textsubscript{2} ≈50\%＋CH\textsubscript{4} ≈30\%）、\ja{コールタール}（\ja{ベンゼン・ナフタレン・フェノール}）。
\item{} \ja{日本はほぼ全量輸入}（\ja{オーストラリア} 6〜7 割）；\ja{国際貿易は生産量の} 1 割強。発電効率：超々臨界 USC ≈42\%、IGCC ≈46〜50\%。\ja{流動層燃焼は} ≈850 ℃ \ja{で} NOx \ja{少・炉内脱硫}。
\end{itemize}
\secthead{3.6 \jas{新燃料・水素}}{\#229\,\#246\,\#285\,\#266\,\#286\,\#40\,\#41\,\#18\,\#75}
\subsecthead{水素（R4-Q12 失分）}{\#229\,\#18}
\begin{gtable}{cYY}
\rowcolor{thead}\textbf{区分} & \textbf{方式} & \textbf{要點} \\
{}製造 & 水蒸気改質＋\ja{シフト反応}（\ja{工業の主流}） & CH\textsubscript{4} ＋ 2H\textsubscript{2}O → CO\textsubscript{2} ＋ 4H\textsubscript{2} → \ja{化石由来は必ず} CO\textsubscript{2}（\ja{グレー}）、CCS 併用＝\ja{ブルー}、\ja{再エネ電解}＝\ja{グリーン} \\
{}輸送 & 圧縮水素 & 19.6 MPa \ja{ボンベ}（\ja{カードル}） \\
{}輸送 & 液化水素 & −253 ℃（≈20 K） \\
{}輸送 & \ja{有機ハイドライド} & \ja{トルエン} ＋ 3H\textsubscript{2} → MCH（\ja{製造地で水素化}＝\hl{発熱}）→ 常温液体輸送 → \ja{需要地で脱水素}（\hl{吸熱}、300〜350 ℃、Pt）→ \ja{トルエン返送} \\
\end{gtable}
\begin{itemize}
\item{} MCH \ja{で体積} ≈1/500。\ja{ベンゼン系を使わない理由}：\ja{発がん性}＋融点 5.5 ℃（\ja{シクロヘキサン} 6.5 ℃）\ja{で冬に凍る}；\ja{トルエン} −95 ℃\ja{・}MCH −127 ℃ \ja{で通年液体}。
\item{} \ja{水素化の発熱}（\ja{結合エネルギー法}）：切断 H–H 436＋π 結合 ≈266 ＝ 702、生成 2×C–H 413 ＝ 826 → ΔH ≈ −124 kJ/mol（C–H \ja{の係数} 2 \ja{を忘れない}；R6-Q14）。
\end{itemize}
\subsecthead{SAF\jas{・新燃料}（R5-Q14）}{\#246\,\#285\,\#40\,\#41\,\#75}
\begin{gtable}{lYY}
\rowcolor{thead}\textbf{略称} & \textbf{原料} & \textbf{工程} \\
{}FT-SPK & \ja{木質バイオマス・都市ごみ}（\ja{雑多な固体}） & \ja{ガス化} → CO＋H\textsubscript{2} → \hl{\ja{フィッシャー・トロプシュ合成}} → 水素化 \\
{}HEFA-SPK & \ja{廃食用油・動植物油脂} & 水素化脱酸素 → 異性化（\ja{最も商業化}） \\
{}HFS-SIP & 糖 & 発酵 → \ja{ファルネセン} → 水素化 \\
{}ATJ-SPK & \ja{エタノール・イソブタノール} & 脱水 → \ja{オリゴマー化} → 水素化 \\
\end{gtable}
FCC \ja{は} SAF \ja{の工程ではない}（\ja{定番の} \batsu{}）。

\begin{itemize}
\item{} \ja{バイオ} ETBE ＝ \ja{バイオエタノール} ＋ \hl{\ja{イソブテン}}（C=C \ja{が必要}；n-\ja{ブタン} \batsu{}）。BDF ＝ 植物油 ＋ \ja{メタノール} \ja{のエステル交換} → \ja{脂肪酸メチルエステル}（FAME）＋ \ja{グリセリン}。DME \ja{は多様な原料から合成ガス経由}。
\item{} 家庭用燃料電池：PEFC ≈70〜80 ℃、SOFC ≈700〜750 ℃（\ja{温度の入れ替えが罠}）。\ja{前段で付臭剤の硫黄を脱硫}。
\item{} 非在来型資源：\ja{オイルシェール}（\ja{未熟成のケロジェン} → 乾留）、\ja{オイルサンド}（\ja{重質ビチューメン} → 蒸気注入）、\ja{タイトオイル・シェールガス}（低浸透率岩 → 水平井＋水圧破砕）、\ja{メタンハイドレート}。「\ja{流動性に優れる}」\batsu{}（\ja{定義と矛盾}）。\ja{原油はオレフィンを含まない}（\ja{分解で生じる}）。
\end{itemize}
\label{ch03-end}
